% ======================================================================
% Chapter 1: Invitation: Pair Production in e+ e- Annihilation
% (book pages 3-12; PDF pages 26-35)
% ======================================================================
\chapter{Invitation: Pair Production in $e^{+}e^{-}$ Annihilation}

The main purpose of Part~I of this book is to develop the basic
calculational method of quantum field theory, the formalism of Feynman
diagrams.  We will then apply this formalism to computations in Quantum
Electrodynamics, the quantum theory of electrons and photons.

Quantum Electrodynamics (QED) is perhaps the best fundamental physical
theory we have.  The theory is formulated as a set of simple equations
(Maxwell's equations and the Dirac equation) whose form is essentially
determined by relativistic invariance.  The quantum-mechanical solutions
of these equations give detailed predictions of electromagnetic phenomena
from macroscopic distances down to regions several hundred times smaller
than the proton.  Feynman diagrams provide for this elegant theory an
equally elegant procedure for calculation:  Imagine a process that can be
carried out by electrons and photons, draw a diagram, and then use the
diagram to write the mathematical form of the quantum-mechanical
amplitude for that process to occur.

In this first part of the book we will develop both the theory of QED and
the method of Feynman diagrams from the basic principles of quantum
mechanics and relativity.  Eventually, we will arrive at a point where we
can calculate observable quantities that are of great interest in the
study of elementary particles.  But to reach our goal of deriving this
simple calculational method, we must first, unfortunately, make a serious
detour into formalism.  The three chapters that follow this one are
almost completely formal, and the reader might wonder, in the course of
this development, where we are going.  We would like to partially answer
that question in advance by discussing the physics of an especially
simple QED process---one sufficiently simple that many of its features
follow directly from physical intuition.  Of course, this intuitive,
bottom-up approach will contain many gaps.  In Chapter~5 we will return
to this process with the full power of the Feynman diagram formalism.
Working from the top down, we will then see all of these difficulties
swept away.

\section{The Simplest Situation}

Since most particle physics experiments involve scattering, the most
commonly calculated quantities in quantum field theory are scattering
cross sections.  We will now calculate the cross section for the simplest
of all QED processes: the annihilation of an electron with its
antiparticle, a positron, to form a pair of heavier leptons (such as
muons).  The existence of antiparticles is actually a prediction of
quantum field theory, as we will discuss in Chapters~2 and~3.  For the
moment, though, we take their existence as given.

An experiment to measure this annihilation probability would proceed by
firing a beam of electrons at a beam of positrons.  The measurable
quantity is the cross section for the reaction
$e^{+}e^{-}\to\mu^{+}\mu^{-}$ as a function of the center-of-mass energy
and the relative angle $\theta$ between the incoming electrons and the
outgoing muons.  The process is illustrated in Fig.~\ref{fig:1.1}.  For
simplicity, we work in the center-of-mass (CM) frame where the momenta
satisfy $\mathbf{p}'=-\mathbf{p}$ and $\mathbf{k}'=-\mathbf{k}$.  We also
assume that the beam energy $E$ is much greater than either the electron
or the muon mass, so that
$|\mathbf{p}|=|\mathbf{p}'|=|\mathbf{k}|=|\mathbf{k}'|=E=E_{\mathrm{cm}}/2$.
(We use boldface type to denote 3-vectors and ordinary italic type to
denote 4-vectors.)

\begin{figure}[t]
  \centering
  \includegraphics[width=0.7\textwidth]{fig1.1.png}
  \caption{The annihilation reaction $e^{+}e^{-}\to\mu^{+}\mu^{-}$, shown
    in the center-of-mass frame.}
  \label{fig:1.1}
\end{figure}

Since both the electron and the muon have spin $1/2$, we must specify
their spin orientations.  It is useful to take the axis that defines the
spin quantization of each particle to be in the direction of its motion;
each particle can then have its spin polarized parallel or antiparallel to
this axis.  In practice, electron and positron beams are often
unpolarized, and muon detectors are normally blind to the muon
polarization.  Hence we should average the cross section over electron and
positron spin orientations, and sum the cross section over muon spin
orientations.

For any given set of spin orientations, it is conventional to write the
differential cross section for our process, with the $\mu^{-}$ produced
into a solid angle $d\Omega$, as
\begin{equation}
  \frac{d\sigma}{d\Omega} = \frac{|\mathcal{M}|^2}{64\pi^2 E_{\mathrm{cm}}^2}.
  \label{eq:1.1}
\end{equation}
The factor $E^{-2}$ provides the correct dimensions for a cross section,
since in our units $(\mathrm{energy})^{-2}\sim(\mathrm{length})^2$.  The
quantity $\mathcal{M}$ is therefore dimensionless; it is the
quantum-mechanical amplitude for the process to occur (analogous to the
scattering amplitude $f$ in nonrelativistic quantum mechanics), and we
must now address the question of how to compute it from fundamental
theory.  The other factors in the expression are purely a matter of
convention.  Equation~\eqref{eq:1.1} is actually a special case, valid for
CM scattering when the final state contains two massless particles, of a
more general formula (whose form cannot be deduced from dimensional
analysis) which we will derive in Section~4.5.

Now comes some bad news and some good news.

The bad news is that even for this simplest of QED processes, the exact
expression for $\mathcal{M}$ is not known.  Actually this fact should come
as no surprise, since even in nonrelativistic quantum mechanics,
scattering problems can rarely be solved exactly.  The best we can do is
obtain a formal expression for $\mathcal{M}$ as a perturbation series in
the strength of the electromagnetic interaction, and evaluate the first
few terms in this series.

The good news is that Feynman has invented a beautiful way to organize
and visualize the perturbation series: the method of Feynman diagrams.
Roughly speaking, the diagrams display the flow of electrons and photons
during the scattering process.  For our particular calculation, the
lowest-order term in the perturbation series can be represented by a
single diagram, shown in Fig.~\ref{fig:1.2}.  The diagram is made up of
three types of components: external lines (representing the four incoming
and outgoing particles), internal lines (representing ``virtual''
particles, in this case one virtual photon), and vertices.  It is
conventional to use straight lines for fermions and wavy lines for
photons.  The arrows on the straight lines denote the direction of
negative charge flow, not momentum.  We assign a 4-momentum vector to
each external line, as shown.  In this diagram, the momentum $q$ of the
one internal line is determined by momentum conservation at either of the
vertices: $q=p+p'=k+k'$.  We must also associate a spin state (either
``up'' or ``down'') with each external fermion.

\begin{figure}[t]
  \centering
  \includegraphics[width=0.55\textwidth]{fig1.2.png}
  \caption{Feynman diagram for the lowest-order term in the
    $e^{+}e^{-}\to\mu^{+}\mu^{-}$ cross section.  At this order the only
    possible intermediate state is a photon ($\gamma$).}
  \label{fig:1.2}
\end{figure}

According to the Feynman rules, each diagram can be translated directly
into a contribution to $\mathcal{M}$.  The rules assign a short algebraic
factor to each element of a diagram, and the product of these factors
gives the value of the corresponding term in the perturbation series.
Getting the resulting expression for $\mathcal{M}$ into a form that is
usable, however, can still be nontrivial.  We will develop much useful
technology for doing such calculations in subsequent chapters.  But we do
not have that technology yet, so to get an answer to our particular
problem we will use some heuristic arguments instead of the actual
Feynman rules.

Recall that in quantum-mechanical perturbation theory, a transition
amplitude can be computed, to first order, as an expression of the form
\begin{equation}
  \mathcal{M} \sim \langle \mathrm{final\ state}\ |\ H_I\ |\ \mathrm{initial\ state}\rangle.
  \label{eq:1.2}
\end{equation}
where $H_I$ is the ``interaction'' part of the Hamiltonian.  In our case
the initial state is $|e^{+}e^{-}\rangle$ and the final state is
$|\mu^{+}\mu^{-}\rangle$.  But our interaction Hamiltonian couples
electrons to muons only through the electromagnetic field (that is,
photons), not directly.  So the first-order result \eqref{eq:1.2}
vanishes, and we must go to the second-order expression
\begin{equation}
  \mathcal{M} \sim
  \sum_{\mu} \langle \mu^{+}\mu^{-}\ |\ H_I\ |\ \gamma\rangle_\mu\;
  \frac{1}{q^2}\;
  \langle \gamma\ |\ H_I\ |\ e^{+}e^{-}\rangle^\mu .
  \label{eq:1.3}
\end{equation}
This is a heuristic way of writing the contribution to $\mathcal{M}$ from
the diagram in Fig.~\ref{fig:1.2}.  The external electron lines correspond
to the factor $|e^{+}e^{-}\rangle$; the external muon lines correspond to
$\langle\mu^{+}\mu^{-}|$.  The vertices correspond to $H_I$, and the
internal photon line corresponds to the operator $|\gamma\rangle\langle
\gamma|$.  We have added vector indices ($\mu$) because the photon is a
vector particle with four components.  There are four possible
intermediate states, one for each component, and according to the rules
of perturbation theory we must sum over intermediate states.  Note that
since the sum in \eqref{eq:1.3} takes the form of a 4-vector dot product,
the amplitude $\mathcal{M}$ will be a Lorentz-invariant scalar as long as
each half of \eqref{eq:1.3} is a 4-vector.

Let us try to guess the form of the vector
$\langle\gamma|\ H_I\ |e^{+}e^{-}\rangle^{\mu}$.  Since $H_I$ couples
electrons to photons with a strength $e$ (the electron charge), the matrix
element should be proportional to $e$.  Now consider one particular set
of initial and final spin orientations, shown in Fig.~\ref{fig:1.3}.  The
electron and muon have spins parallel to their directions of motion; they
are ``right-handed''.  The antiparticles, similarly, are ``left-handed''.
The electron and positron spins add up to one unit of angular momentum in
the $+z$ direction.  Since $H_I$ should conserve angular momentum, the
photon to which these particles couple must have the correct
polarization vector to give it this same angular momentum:
\begin{equation}
  \epsilon^\mu = (0, 1, i, 0).
\end{equation}
Thus we have
\begin{equation}
  \langle\gamma|\ H_I\ |e^{+}e^{-}\rangle^\mu
    \propto e\,(0, 1, i, 0).
  \label{eq:1.4}
\end{equation}
The muon matrix element should, similarly, have a polarization
corresponding to one unit of angular momentum along the direction of the
$\mu^{-}$ momentum $\mathbf{k}$.  To obtain the correct vector, rotate
\eqref{eq:1.4} through an angle $\theta$ in the $xz$-plane:
\begin{equation}
  \langle\gamma|\ H_I\ |\mu^{+}\mu^{-}\rangle^\mu
    \propto e\,(0,\ \cos\theta,\ i,\ -\sin\theta).
  \label{eq:1.5}
\end{equation}
To compute the amplitude $\mathcal{M}$, we complex-conjugate this vector
and dot it into \eqref{eq:1.4}.  Thus we find, for this set of spin
orientations,
\begin{equation}
  \mathcal{M}(RL \to RL) = -e^{2}\,(1 + \cos\theta).
  \label{eq:1.6}
\end{equation}
Of course we cannot determine the overall factor by this method, but in
\eqref{eq:1.6} it happens to be correct, thanks to the conventions adopted
in \eqref{eq:1.1}.  Note that the amplitude vanishes for
$\theta=180^{\circ}$, just as one would expect: A state whose angular
momentum is in the $+z$ direction has no overlap with a state whose
angular momentum is in the $-z$ direction.

\begin{figure}[t]
  \centering
  \includegraphics[width=0.6\textwidth]{fig1.3.png}
  \caption{One possible set of spin orientations.  The electron and the
    negative muon are right-handed, while the positron and the positive
    muon are left-handed.}
  \label{fig:1.3}
\end{figure}

Next consider the case in which the electron and positron are both
right-handed.  Now their total spin angular momentum is zero, and the
argument is more subtle.  We might expect to obtain a longitudinally
polarized photon with a Clebsch--Gordan coefficient of $1/\sqrt{2}$, just
as when we add angular momenta in three dimensions,
$|\tfrac{1}{2}\tfrac{1}{2}\rangle =
(1/\sqrt{2})(|j{=}1, m{=}0\rangle + |j{=}0, m{=}0\rangle)$.  But we are
really adding angular momenta in the four-dimensional Lorentz group, so
we must take into account not only spin (the transformation properties of
states under rotations), but also the transformation properties of states
under boosts.  It turns out, as we shall discuss in Chapter~3, that the
Clebsch--Gordan coefficient that couples a 4-vector to the state
$|e^{+}e^{-}\rangle$ of massless fermions is zero.  (For the record, the
state is a superposition of scalar and antisymmetric tensor pieces.)  Thus
the amplitude $\mathcal{M}(RR\to RL)$ is zero, as are the eleven other
amplitudes in which either the initial or final state has zero total
angular momentum.

The remaining nonzero amplitudes can be found in the same way that we
found the first one.  They are
\begin{align}
  \mathcal{M}(RL \to LR) &= -e^{2}\,(1 - \cos\theta),\\
  \mathcal{M}(LR \to RL) &= -e^{2}\,(1 - \cos\theta),\\
  \mathcal{M}(LR \to LR) &= -e^{2}\,(1 + \cos\theta).
  \label{eq:1.7}
\end{align}
Inserting these expressions into \eqref{eq:1.1}, averaging over the four
initial-state spin orientations, and summing over the four final-state
spin orientations, we find
\begin{equation}
  \frac{d\sigma}{d\Omega} =
    \frac{\alpha^{2}}{4E_{\mathrm{cm}}^{2}}(1 + \cos^{2}\theta),
  \label{eq:1.8}
\end{equation}
where $\alpha = e^{2}/4\pi \approx 1/137$.  Integrating over the angular
variables $\theta$ and $\phi$ gives the total cross section,
\begin{equation}
  \sigma_{\mathrm{total}} =
    \frac{4\pi\alpha^{2}}{3E_{\mathrm{cm}}^{2}} .
  \label{eq:1.9}
\end{equation}
Results \eqref{eq:1.8} and \eqref{eq:1.9} agree with experiments to about
$10\%$; almost all of the discrepancy is accounted for by the next term in
the perturbation series, corresponding to the diagrams shown in
Fig.~\ref{fig:1.4}.  The qualitative features of these expressions---the
angular dependence and the sharp decrease with energy---are obvious in the
actual data.  (The properties of these results are discussed in detail in
Section~5.1.)

\begin{figure}[t]
  \centering
  \includegraphics[width=0.8\textwidth]{fig1.4.png}
  \caption{Feynman diagrams that contribute to the $\alpha^{3}$ term in
    the $e^{+}e^{-}\to\mu^{+}\mu^{-}$ cross section.}
  \label{fig:1.4}
\end{figure}

\section{Embellishments and Questions}

We obtained the angular distribution predicted by Quantum Electrodynamics
for the reaction $e^{+}e^{-}\to\mu^{+}\mu^{-}$ by applying angular
momentum arguments, with little appeal to the underlying formalism.
However, we used the simplifying features of the high-energy limit and the
center-of-mass frame in a very strong way.  The analysis we have presented
will break down when we relax any of our simplifying assumptions.  So how
does one perform general QED calculations?  To answer that question we
must return to the Feynman rules.

As mentioned above, the Feynman rules tell us to draw the diagram(s) for
the process we are considering, and to associate a short algebraic factor
with each piece of each diagram.  Figure~\ref{fig:1.5} shows the diagram
for our reaction, with the various assignments indicated.

\begin{figure}[t]
  \centering
  \includegraphics[width=0.55\textwidth]{fig1.5.png}
  \caption{Diagram of Fig.~\ref{fig:1.2}, with expressions corresponding
    to each vertex, internal line, and external line.}
  \label{fig:1.5}
\end{figure}

For the internal photon line we write $-ig_{\mu\nu}/q^{2}$, where
$g_{\mu\nu}$ is the usual Minkowski metric tensor and $q$ is the 4-momentum
of the virtual photon.  This factor corresponds to the operator
$|\gamma\rangle\langle\gamma|$ in our heuristic expression \eqref{eq:1.3}.
For each vertex we write $-ie\gamma^{\mu}$, corresponding to $H_I$ in
\eqref{eq:1.3}.  The objects $\gamma^{\mu}$ are a set of four $4\times4$
constant matrices.  They do the ``addition of angular momentum'' for us,
coupling a state of two spin-$1/2$ particles to a vector particle.

The external lines carry expressions for four-component column-spinors
$u$, $v$, or row-spinors $\bar{u}$, $\bar{v}$.  These are essentially the
momentum-space wavefunctions of the initial and final particles, and
correspond to $|e^{+}e^{-}\rangle$ and $\langle\mu^{+}\mu^{-}|$ in
\eqref{eq:1.3}.  The indices $s$, $s'$, $r$, and $r'$ denote the spin
state, either up or down.

We can now write down an expression for $\mathcal{M}$, reading everything
straight off the diagram:
\begin{equation}
  \mathcal{M} =
    \bar{v}^{s'}(p')(-ie\gamma^{\mu})u^{s}(p)
    \frac{-ig_{\mu\nu}}{q^{2}}
    \bar{u}^{r}(k)(-ie\gamma^{\nu})v^{r'}(k')
    = \frac{(-ie)^{2}}{q^{2}}
      \Bigl(\bar{v}^{s'}(p')\gamma^{\mu}u^{s}(p)\Bigr)
      \Bigl(\bar{u}^{r}(k)\gamma_{\mu}v^{r'}(k')\Bigr).
  \label{eq:1.10}
\end{equation}
It is instructive to compare this in detail with Eq.~\eqref{eq:1.3}.

To derive the cross section \eqref{eq:1.8} from \eqref{eq:1.10}, we could
return to the angular momentum arguments used above, supplemented with
some concrete knowledge about $\gamma$ matrices and Dirac spinors.  We will
do the calculation in this manner in Section~5.2.  There are, however, a
number of useful tricks that can be employed to manipulate expressions
like \eqref{eq:1.10}, especially when one wants to compute only the
unpolarized cross section.  Using this ``Feynman trace technology''
(so-called because one must evaluate traces of products of
$\gamma$-matrices), it isn't even necessary to have explicit expressions
for the $\gamma$-matrices and Dirac spinors.  The calculation becomes
almost completely mindless, and the answer \eqref{eq:1.8} is obtained
after less than a page of algebra.

But since the Feynman rules and trace technology are so powerful, we can
also relax some of our simplifying assumptions.  To conclude this section,
let us discuss several ways in which our calculation could have been more
difficult.

The easiest restriction to relax is that the muons be massless.  If the
beam energy is not much greater than the mass of the muon, all of our
predictions should depend on the ratio $m_\mu/E_{\mathrm{cm}}$.  (Since the
electron is 200 times lighter than the muon, it can be considered massless
whenever the beam energy is large enough to create muons.)  Using Feynman
trace technology, it is extremely easy to restore the muon mass to our
calculation.  The amount of algebra is increased by about fifty percent,
and the relation \eqref{eq:1.1} between the amplitude and the cross
section must be modified slightly, but the answer is worth the effort.
We do this calculation in detail in Section~5.1.

Working in a different reference frame is also easy; the only modification
is in the relation \eqref{eq:1.1} between the amplitude and the cross
section.  Or one can simply perform a Lorentz transformation on the CM
result, boosting it to a different frame.

When the spin states of the initial and/or final particles are known and
we still wish to retain the muon mass, the calculation becomes somewhat
cumbersome but no more difficult in principle.  The trace technology can
be generalized to this case, but it is often easier to evaluate expression
\eqref{eq:1.10} directly, using the explicit values of the spinors $u$ and
$v$.

Next one could compute cross sections for different processes.  The
process $e^{+}e^{-}\to e^{+}e^{-}$, known as Bhabha scattering, is more
difficult because there is a second allowed diagram (see
Fig.~\ref{fig:1.6}).  The amplitudes for the two diagrams must first be
added, then squared.

\begin{figure}[t]
  \centering
  \includegraphics[width=0.6\textwidth]{fig1.6.png}
  \caption{The two lowest-order diagrams for Bhabha scattering,
    $e^{+}e^{-}\to e^{+}e^{-}$.}
  \label{fig:1.6}
\end{figure}

Other processes contain photons in the initial and/or final states.  The
paradigm example is Compton scattering, for which the two lowest-order
diagrams are shown in Fig.~\ref{fig:1.7}.  The Feynman rules for external
photon lines and for internal electron lines are no more complicated than
those we have already seen.  We discuss Compton scattering in detail in
Section~5.5.

\begin{figure}[t]
  \centering
  \includegraphics[width=0.6\textwidth]{fig1.7.png}
  \caption{The two lowest-order diagrams for Compton scattering.}
  \label{fig:1.7}
\end{figure}

Finally we could compute higher-order terms in the perturbation series.
Thanks to Feynman, the diagrams are at least easy to draw; we have seen
those that contribute to the next term in the $e^{+}e^{-}\to\mu^{+}\mu^{-}$
cross section in Fig.~\ref{fig:1.4}.  Remarkably, the algorithm that
assigns algebraic factors to pieces of the diagrams holds for all
higher-order contributions, and allows one to evaluate such diagrams in a
straightforward, if tedious, way.  The computation of the full set of nine
diagrams is a serious chore, at the level of a research paper.

In this book, starting in Chapter~6, we will analyze much of the physics
that arises from higher-order Feynman diagrams such as those in
Fig.~\ref{fig:1.4}.  We will see that the last four of these diagrams,
which involve an additional photon in the final state, are necessary
because no detector is sensitive enough to notice the presence of
extremely low-energy photons.  Thus a final state containing such a photon
cannot be distinguished from our desired final state of just a muon pair.

The other five diagrams in Fig.~\ref{fig:1.4} involve intermediate states
of several virtual particles rather than just a single virtual photon.
In each of these diagrams there will be one virtual particle whose
momentum is not determined by conservation of momentum at the vertices.
Since perturbation theory requires us to sum over all possible
intermediate states, we must integrate over all possible values of this
momentum.  At this step, however, a new difficulty appears: The
loop-momentum integrals in the first three diagrams, when performed
naively, turn out to be infinite.  We will provide a fix for this problem,
so that we get finite results, by the end of Part~I.  But the question of
the physical origin of these divergences cannot be dismissed so lightly;
that will be the main subject of Part~II of this book.

We have discussed Feynman diagrams as an algorithm for performing
computations.  The chapters that follow should amply illustrate the power
of this tool.  As we expose more applications of the diagrams, though,
they begin to take on a life and significance of their own.  They indicate
unsuspected relations between different physical processes, and they
suggest intuitive arguments that might later be verified by calculation.
We hope that this book will enable you, the reader, to take up this tool
and apply it in novel and enlightening ways.
